\section{Ehrhart Theory and Lattice Polytopes}

\newcommand{\Z}{\mathbb{Z}}
\newcommand{\Q}{\mathbb{Q}}
\newcommand{\R}{\mathbb{R}}


Group members: Anouk Brose, Jesús De Loera, Noah Gießing, Mei Han, Santiago Morales Duarte, Justus Springer
\bigskip

Our group set out with the long-term goal to formalize Ehrhart's Theorem, or some version of it, in mathlib. We took this as an overarching goal to work towards, and used these two weeks to contribute some first lemmas, learn Lean, both in writing and in mathlib conventions, as well as (for some of us) to get experience with GitHub. Our first steps towards formalizing Ehrhart's Theorem were to define a suitable notion of \textit{lattice polytopes} in Lean, improve the (algebraic) definition of \textit{lattice} (like $\mathbb{Z}^d$) and to prove a commutative algebra version of \textit{Gordan's lemma}.
 \bigskip

\textit{Lattice Polytopes:} We discussed if defining a lattice polytope should be a property of \lstinline{IsPolytope} or a definition on top of the bundled structure \lstinline|Polytope|. Currently \lstinline{IsPolytope} is defined as `there exists a set S whose convex hull is the polytope'. The definition for a lattice polytope should be similar, only with the condition that the set S should be a subset of the lattice. We quickly decided that this should not be a definition on top of either one of the polytope definitions, because then we have two existential quantifiers describing the object. For now we have defined a version, called \lstinline|IsGroundedPolytope| in parallel to the definition \lstinline|IsPolytope|, with the extra condition that the generating set is from the lattice. We used the word \textit{grounded} for now, as in that this definition can work for any `ground set', not just a lattice. Later it will be of interest to have definitions for \textit{rational polytopes}, and in literature there's also a notion of \textit{S-polytopes} where the vertices come from a specific subset $S$, e.g., integer vectors with prime coordinates. Another option would be to not define a new notion at all, but instead, whenever a lattice polytope is used, just write a polytope with extreme points in the lattice (or the grounded set).
\medskip

Our work is located in the Polyhedral library under \path{Geometry/Convex/ConvexSpace/Polytope}. It builds on the definition of polytope.
\medskip

\textit{Lattices:} There are two definitions that define a lattice (such as $\Z^d$) in mathlib, but we don't think they give the correct framework of lattices. First of all both definitions only define full-dimensional lattices (w.r.t. their ambient space). One definition \lstinline|IsZLattice| is only defined in normed spaces, and the other definition \lstinline|Submodule.IsLattice| is an algebraic notion, but it is designed for an a $R$-lattice with a $\text{Frac}(R)$ vector space in mind. 
For example, the current definition lets the $\Z$-span of $1,\sqrt{2}$ in $\R$ be a $\Z$-lattice (because, in this case, it was designed with $\Q$-vector spaces and not $\R$-vector spaces). We started working on a definition, \lstinline|Submodule.IsLattice'|, that should replace \lstinline|Submodule.IsLattice|, but allows for lower-dimensional lattices, and works well in ambient spaces that are not $\text{Frac}(R)$-vector spaces. 

\medskip
Our work is in Mathlib's \path{Algebra/Module/Submodule/Basic.lean} folder, because that is where the current \lstinline|Submodule.IsLattice| is located. This project was independent of the polyhedral library.
\medskip 

\textit{Gordan's Lemma:} This is a crucial lemma to prove Ehrhart's theorem for simplices. It says that the lattice points in a rational simplicial cone admit a unique decomposition as a sum of a lattice point in the fundamental parallelepiped defined by integer rays of the cone, and a non-negative integer combination of the rays. We discussed for a while in which generality to formalize this lemma. A more basic version would say that an additive monoid in any cone generated by rays from the monoid is finitely generated. However, to prove Ehrhart's theorem we will need the exact decomposition theorem, so we formalized it in that setting. We get the finite generation as a corollary.
\medskip

Our work is in Mathlib's \path{Geometry/Convex/Cone/Pointed folder}. This project ended up being independent of the polyhedral library.
\medskip 

\textbf{Mathlib Contributions:} We successfully contributed two small pull requests to Mathlib. (\#43107 and \#43297.)
The first was a correction in the file \lstinline|HilbertPoly|. We updated the lemma, now called \lstinline|preHilbertPoly_eq_choose_add_sub|, which says that the polynomial $\binom{n-k+d}{d}$ evaluates to the corresponding binomial coefficient. Originally this lemma was only stated under the condition that $k \leq n$, but in fact this holds for all $k \leq n+d$. This formulation will be useful in later steps to prove Ehrhart's theorem. The second contribution was on a lemma about the rank of submodules, which we noticed while defining \lstinline|IsLattice'|, and will be useful to prove basic lemmas about our new definition of lattices.
\medskip 

\textbf{Next Steps.} We plan on continuing the project and our next steps are to (1) finish developing the definition \lstinline|IsLattice'|, (2) connect Gordan's lemma to prove identities on generating functions, both on the level of multivariate generating functions and univariate generating functions. 
